% Part: incompleteness
% Chapter: introduction
% Section: overview

\documentclass[../../../include/open-logic-section]{subfiles}

\begin{document}

\olfileid{inc}{int}{ovr}

\olsection{অসম্পূর্ণতা-ফলগুলির সারসংক্ষেপ}

হিলবার্ট আশা করেছিলেন, গণিতকে এমন একটি !!{axiomatizable} তত্ত্বে বিধিবদ্ধ করা
যাবে, যেটি !!{complete} এবং !!{decidable}—এ কথা প্রমাণ করাও সম্ভব হবে।
এ ছাড়া, তিনি খুব দুর্বল, ``সসীম-পদ্ধতিগত'' উপায়ে এই তত্ত্বের সঙ্গতি প্রমাণ
করতে চেয়েছিলেন; এর ফলে অন্তর্দৃষ্টিবাদের চ্যালেঞ্জের বিরুদ্ধে ধ্রুপদী গণিতকে
রক্ষা করা যেত। গ্যোডেলের অসম্পূর্ণতা উপপাদ্য দেখাল যে এই লক্ষ্যগুলি অর্জন করা
যায় না।

গ্যোডেলের প্রথম অসম্পূর্ণতা উপপাদ্য দেখিয়েছিল যে রাসেল ও হোয়াইটহেডের
\emph{Principia Mathematica}-র একটি সংস্করণ !!{complete} নয়। কিন্তু প্রমাণটি
আসলে অত্যন্ত সাধারণ এবং বহু ধরনের তত্ত্বে প্রযোজ্য। অর্থাৎ শুধু
\emph{Principia Mathematica}-ই গণিতকে সম্পূর্ণভাবে ধারণ করতে পারেনি এমন নয়,
কোনো গ্রহণযোগ্য তত্ত্বই তা পারে না। অসম্পূর্ণতা উপপাদ্য প্রয়োগের জন্য যথেষ্ট
তত্ত্বগত বৈশিষ্ট্যগুলি আলাদা করে শনাক্ত করতে এবং গ্যোডেলের প্রমাণকে সেই
বৈশিষ্ট্যগুলির উপরই নির্ভরশীল করে সাধারণীকরণ করতে কিছু সময় লেগেছিল। এখন আমরা
পাটিগণিতের ভাষা $\Lang L_A$-তে তত্ত্বগুলির জন্য প্রথম অসম্পূর্ণতা উপপাদ্যের
একটি অত্যন্ত সাধারণ রূপ বলতে পারি।

\begin{thm}
যদি $\Gamma$ একটি সঙ্গতিপূর্ণ ও !!{axiomatizable} তত্ত্ব হয়, যা
$\Lang L_A$-এর মধ্যে সব গণনসাধ্য অপেক্ষক ও !!{decidable} সম্বন্ধকে
!!{represents} করে, তবে $\Gamma$ !!{complete} নয়।
\end{thm}

$\Gamma$ !!{complete} নয় বলা মানে অন্তত একটি !!{sentence}~$!A$-এর জন্য
$\Gamma \Proves/ !A$ এবং $\Gamma \Proves/ \lnot !A$। এই রকম
!!a{sentence}-কে $\Gamma)$-এর সাপেক্ষে \emph{স্বাধীন} বলা হয়। সত্যিই খুব
সহজেই প্রমাণ করা যায় যে স্বাধীন বাক্য থাকবেই। কিন্তু গ্যোডেলের
উপপাদ্যের শক্তি হলো, এটি এমন একটি স্বাধীন !!{sentence}-এর \emph{নির্দিষ্ট
উদাহরণ} দেয়। এই চমৎকার নির্মাণে !!a{sentence}~$!G_\Gamma$ পাওয়া যায়,
যাকে $\Gamma$-এর জন্য \emph{গ্যোডেল বাক্য} বলা হয়; এটি অপ্রমাণযোগ্য, কারণ
$\Gamma$-তে $!G_\Gamma$ সেই দাবির সমতুল্য যে $!G_\Gamma$-কে $\Gamma$-তে
প্রমাণ করা যায় না। নির্মাণটি \emph{গঠনমূলক}: অর্থাৎ $\Gamma$-এর
স্বতঃসিদ্ধায়ন এবং !!{derivation} ব্যবস্থার বর্ণনা দেওয়া থাকলে, প্রমাণটি
$!G_\Gamma$ বাস্তবে লিখে ফেলার একটি পদ্ধতি দেয়।

গ্যোডেলের প্রমাণে এমন একটি পদ্ধতি দরকার যাতে $\Lang L_A$-তেই $\Lang L_A$-এর
নিজস্ব পদ ও !!{formula}s-গুলির ধর্ম ও ক্রিয়া প্রকাশ করা যায়। যেমন,
``$!A$ একটি !!a{sentence}'', ``$\delta$ হলো $!A$-এর !!a{derivation}''—এই
ধর্মগুলি এবং $\Subst{!A}{t}{x}$-এর মতো ক্রিয়াগুলি প্রকাশ করতে হবে। এই
পদ্ধতিকে (ক) চিহ্ন ও চিহ্ন-অনুক্রমকে ``কোড'' করে প্রাকৃতিক সংখ্যায় প্রকাশ
করতে হবে (পদ, !!{formula}s, !!{derivation}s ইত্যাদি আসলে এমন অনুক্রম, আর
$\Lang L_A$ প্রাকৃতিক সংখ্যা নিয়ে কথা বলতে পারে), এবং (খ) এমনভাবে করতে হবে
যাতে $\Gamma$ প্রাসঙ্গিক তথ্য প্রমাণ করতে পারে। তাই দেখাতে হবে যে এই ধর্ম ও
সম্বন্ধগুলি প্রাকৃতিক সংখ্যার !!{decidable} ধর্ম দিয়ে কোড হয় এবং ক্রিয়াগুলি
প্রাকৃতিক সংখ্যার উপর গণনসাধ্য অপেক্ষকের অনুরূপ। একে ``syntax-এর
arithmetization'' বলা হয়।

syntax কীভাবে arithmetize করা যায় তা বিশদে দেখার আগে, রূপরেখা হিসেবে আমরা
$\Gamma$ ``যথেষ্ট শক্তিশালী''—অর্থাৎ সব গণনসাধ্য অপেক্ষক ও !!{decidable}
সম্বন্ধকে !!{represents} করে—এই শর্তটি বিবেচনা করব। এর জন্য ``গণনসাধ্য''-র
একটি সুনির্দিষ্ট সংজ্ঞা দরকার। এটি নানা ভাবে দেওয়া যায়, যেমন টিউরিং
মেশিনের মডেল দিয়ে, অথবা কোনো সাধারণ-উদ্দেশ্য প্রোগ্রামিং ভাষায় যে অপেক্ষকগুলি
গণনা করা যায় সেগুলি দিয়ে। কিন্তু $\Lang L_A$-এর তত্ত্বে এই অপেক্ষক ও
সম্বন্ধগুলি !!{represents}s করাই আমাদের লক্ষ্য বলে, শুধু সংখ্যাগত অপেক্ষকের
একটি সরল সংজ্ঞা নেওয়াই ভালো। এটিই \emph{recursive function}-এর ধারণা। তাই
পরবর্তী আলোচনায় recursive function আগে আসবে; তারপর দেখানো হবে যে
$\Th{Q}$ ইতিমধ্যেই সব recursive function ও সম্বন্ধকে !!{represents} করে।
তখন $\Th{Q}$ ও $\Th{PA}$-র মতো নির্দিষ্ট তত্ত্বে অসম্পূর্ণতা উপপাদ্য প্রয়োগ
করা যাবে, কারণ তারা ``যথেষ্ট শক্তিশালী'' তত্ত্বের উদাহরণ হবে।
% BN-SRC-201 TARGET-CORRECTION-BEGIN
\par\smallskip\noindent\textnormal{\small\textbf{উৎস-সংশোধন (BN-SRC-201):} হিমায়িত রূপরেখায় ``আগে'' ও ``তারপর'' কথাগুলি বাস্তব অধ্যায়-আমদানি-ক্রমের সঙ্গে মেলে না—অঙ্কিতকরণ-সংক্রান্ত অধ্যায়টি প্রতিনিধিত্বের অধ্যায়ের আগে এসেছে। বাংলা পাঠে এই বাক্যগুলিকে রূপরেখার যৌক্তিক নির্ভরতা হিসেবে চিহ্নিত করে ``পরবর্তী আলোচনায়'' বলা হয়েছে, যাতে পাঠক আমদানি-ক্রমকে ভুলভাবে না পড়েন।} % BN-SRC-201 TARGET-CORRECTION-END

syntax-এর arithmetization-এর শেষ ফল হলো !!a{formula}~$\OProv[\Gamma](x)$,
যা !!{formula}s-গুলিকে সংখ্যায় কোড করার মাধ্যমে $\Gamma$-এর স্বতঃসিদ্ধগুলি
থেকে প্রমাণযোগ্যতা প্রকাশ করে। নির্দিষ্টভাবে, $!A$-কে যদি $n$ সংখ্যা দিয়ে
কোড করা হয় এবং $\Gamma \Proves !A$ হয়, তবে
$\Gamma \Proves \OProv[\Gamma](\num{n})$। এই ``provability predicate''-এর
সাহায্যে $\Gamma$-এর একটি নির্দিষ্ট অর্থে $\Lang L_A$-এর !!a{sentence} হিসেবে $\Gamma$-র
সঙ্গিতও প্রকাশ করা যায়: $\Gamma$-র ``consistency statement'' হিসেবে
!!{sentence} $\lnot \OProv[\Gamma](\num{n})$ নিই, যেখানে $n$ হলো কোনো
বিরোধের কোড, যেমন~$\lfalse$-এর কোড। দ্বিতীয় অসম্পূর্ণতা উপপাদ্য বলে,
সঙ্গতিপূর্ণ !!{axiomatizable} তত্ত্বগুলি নিজেদের consistency statement-ও
প্রমাণ করতে পারে না। এই উপপাদ্য প্রয়োগের শর্ত শুধু সব গণনসাধ্য অপেক্ষক ও
!!{decidable} সম্বন্ধকে প্রতিনিধিত্ব করার চেয়ে কিছুটা শক্তিশালী; আমরা দেখাব,
$\Th{PA}$ সেই শর্তগুলি পূরণ করে।
% BN-SRC-202 TARGET-CORRECTION-BEGIN
\par\smallskip\noindent\textnormal{\small\textbf{উৎস-সংশোধন (BN-SRC-202):} হিমায়িত অনুচ্ছেদের শেষ অংশে আগে সংজ্ঞায়িত বস্তুভাষার সূত্র $\OProv[\Gamma](x)$-এর বদলে দুবার ধাতুভাষিক $\Prov[\Gamma]$ লেখা হয়েছে। বাংলা পাঠে উভয় সূত্রেই বস্তুভাষার $\OProv[\Gamma]$ পুনঃস্থাপন করা হয়েছে।} % BN-SRC-202 TARGET-CORRECTION-END

\end{document}
